% Generated draft: requires mathematical and source review.
\documentclass[11pt]{article}
\usepackage[margin=25mm]{geometry}
\usepackage{fontspec}
\setmainfont{Latin Modern Roman}
\usepackage{amsmath,amssymb,mathtools,unicode-math}
\setmathfont{Latin Modern Math}
\usepackage{enumitem,needspace,xurl,hyperref}
\setlength{\parindent}{0pt}
\setlength{\parskip}{5pt}
\setlength{\emergencystretch}{3em}
\newcommand{\N}{\mathbb{N}}
\newcommand{\Z}{\mathbb{Z}}
\newcommand{\Q}{\mathbb{Q}}
\newcommand{\R}{\mathbb{R}}
\newcommand{\C}{\mathbb{C}}
\newcommand{\F}{\mathbb{F}}
\newcommand{\A}{\mathbb{A}}
\newcommand{\PP}{\mathbb{P}}
\newcommand{\E}{\mathbb{E}}
\newcommand{\eps}{\varepsilon}
\newcommand{\dd}{\,\mathrm{d}}
\begin{document}
\section*{UnsolvedMath ID: 30001572}
\textbf{Source problem code:} catalogue entry
\subsection{Definitions and mathematical statement}
Let \(X\) be a smooth complex projective Fano manifold of complex dimension \(n\), so that its anticanonical line bundle \(-K_X\) is ample. For a positive integer \(m\), the complete linear system \(|-mK_X|\) parametrizes effective divisors linearly equivalent to \(-mK_X\).

Let \(\phi_\infty\) denote the subsequential limit associated with the normalized Kähler--Ricci flow mentioned in the supplied record. Write \(\mathcal I(\gamma\phi_\infty)\) for the multiplier ideal intended by that record, and, if this ideal is defined and coherent, write
\[
V(\gamma\phi_\infty)
:=\bigl(X,\mathcal O_X/\mathcal I(\gamma\phi_\infty)\bigr)
\]
for the corresponding closed subscheme of \(X\). This notation does not supply the missing definition: the record does not specify the allowed values of \(\gamma\), the normalization and convergence used to obtain \(\phi_\infty\), or the local integrability convention defining the ideal.

The expanded statement supplies the following prescription: for every real \(\varepsilon>0\) and all sufficiently large positive integers \(m\), choose a \(G\)-invariant linear system \(|\Sigma|\subset|-mK_X|\) and a member \(D\in|\Sigma|\) such that
\[
\left(X,c_{\varepsilon,m}D\right)
\quad\text{is not Kawamata log terminal},
\qquad
c_{\varepsilon,m}:=\frac{n+\varepsilon}{(n+1)m}.
\]
Here, for an effective divisor \(D\) and a real number \(c\geq0\), the pair \((X,cD)\) is Kawamata log terminal if
\[
A_X(E)-c\,\operatorname{ord}_E(D)>0
\]
for every prime divisor \(E\) over \(X\), where \(A_X(E)\) is its log discrepancy over \(X\). The group \(G\), its action, and the hypotheses ensuring the prescribed choices exist are not supplied. The prescription is retained as source data, not asserted here as an existence theorem for every Fano manifold. Invariance is stated for the linear system; invariance of its chosen member \(D\) is not stated.

Subject to resolving these omissions, the source asks:
\begin{enumerate}
\item Is \(V(\gamma\phi_\infty)\) a divisor on \(X\)? The record does not specify whether this means an effective Cartier divisor as a closed subscheme, or only a divisorial support or associated cycle.
\item If it is a divisor, how does it differ from the divisor \(D\) referred to as belonging to ``Fact 2'' in the original statement? The expanded statement describes \(D\) by the prescription above, but the text and hypotheses of Fact 2 are absent. No particular comparison relation is specified.
\item Which of \(V(\gamma\phi_\infty)\) and \(D\) violates Ross--Thomas stability of \(X\)? The polarization and precise stability criterion are not specified, so this cannot yet be expressed as a unique numerical instability condition.
\end{enumerate}
The conditional wording of the second question is retained from the original statement. These are questions for review, not assertions that either object is divisorial or destabilizing.
\subsection{Short English statement}
The source asks whether a multiplier-ideal subscheme arising from a normalized Kähler--Ricci flow limit on a Fano manifold is a divisor; if so, how it compares with a divisor described through a non-Kawamata-log-terminal pair; and which object violates Ross--Thomas stability. Missing definitions and hypotheses prevent a fully determined formulation. The supplied open status is a dated curation label, not an independently verified conclusion.
\subsection{Sources}
\begin{itemize}
\item Local UnsolvedMath registry, supplied complete record with permanent ID \(30001572\), source problem code \texttt{catalogue entry}, title ``Multiplier-Ideal Divisors from Kähler--Ricci Flow Limits,'' in the collection ``Oberwolfach Reports Open Problems.'' No individual proposer is identified; the supplied proposed year is \(2010\).
\item Supplied report citation: \emph{Komplexe Analysis (2011)}, Oberwolfach Reports, pp.\ 2283--2333, DOI \texttt{10.4171/owr/2010/39}. Supplied source URL: \url{https://doi.org/10.4171/owr/2010/39}. The background attributes the problem to workshop open-problem discussions documented in Oberwolfach Reports.
\item Supplied literature source, labeled ``original OWR report'': \url{https://doi.org/10.4171/owr/2010/39}.
\item Literature assessment dated exactly \texttt{2026-08-21}: the supplied curation labels the item open and says that the extraction contains three questions but omits definitions of the ideal and comparison divisor. It treats the item as unresolved rather than as a verifiable solved or disproved claim. This is supplied dated curation, not independently verified literature assessment; its status and cited claims require independent verification.
\item The registry also supplies the status \texttt{corrected\_verified} and a statement-verification note claiming comparison with the official report and normalization without changing the source-level claim. Those are supplied provenance claims; this draft does not independently confirm them or claim to have consulted the report.
\end{itemize}
\end{document}
